% ============================================================
% Table 1: Synthetic Option Chain — Calibration Setup
% ============================================================
\begin{table}[htbp]
  \centering
  \caption{Synthetic Option Chain Used for Calibration Experiments}
  \label{tab:synthetic_setup}
  \begin{tabular}{ll}
    \toprule
    \textbf{Setting} & \textbf{Value} \\
    \midrule
    \multicolumn{2}{l}{\textit{Ground Truth Heston Parameters}} \\
    \quad Mean-reversion speed $\kappa$       & 2.0 \\
    \quad Long-run variance $\theta$          & 0.1 \\
    \quad Volatility of variance $\sigma$     & 0.4 \\
    \quad Correlation $\rho$                  & $-0.6$ \\
    \quad Initial variance $v_0$              & 0.05 \\
    \quad Risk-free rate $r$                  & 0.05 \\
    \quad Spot price $S_0$                    & 1000.0 \\
    \quad Feller condition ($2\kappa\theta > \sigma^2$) & Satisfied ($0.40 > 0.16$) \\
    \midrule
    \multicolumn{2}{l}{\textit{Option Chain Grid}} \\
    \quad Maturities $\tau$   & $\{0.3,\, 0.4,\, 0.5,\, 0.6,\, 0.7\}$ (years) \\
    \quad Moneyness $K/S_0$   & 20 levels $\in [0.6,\, 1.5]$ \\
    \quad Valid contracts $M$ & 99 \\
    \midrule
    \multicolumn{2}{l}{\textit{Optimisation (shared by both methods)}} \\
    \quad Algorithm          & Multi-start Adam \\
    \quad Number of restarts & 10 \\
    \quad Learning rate      & $5\times10^{-3}$ \\
    \quad Max iterations     & 500 per restart \\
    \quad Early-stopping patience & 40 steps \\
    \bottomrule
  \end{tabular}
\end{table}


% ============================================================
% Table 2: Parameter Recovery — Proposed vs Zhang (2025)
% ============================================================
\begin{table}[htbp]
  \centering
  \caption{Heston Parameter Recovery on Synthetic Data: Proposed Method vs.\ Zhang~(2025)}
  \label{tab:param_recovery}
  \setlength{\tabcolsep}{6pt}
  \begin{tabular}{lccccc}
    \toprule
    \multirow{2}{*}{\textbf{Parameter}} &
    \multirow{2}{*}{\textbf{True}} &
    \multicolumn{2}{c}{\textbf{Proposed Method}} &
    \multicolumn{2}{c}{\textbf{Zhang (2025)}} \\
    \cmidrule(lr){3-4}\cmidrule(lr){5-6}
    & & \textbf{Calibrated} & \textbf{Rel.\ Err.} & \textbf{Calibrated} & \textbf{Rel.\ Err.} \\
    \midrule
    $\kappa$  & 2.000 & 3.819 & 90.94\% & 3.588 & 79.42\% \\
    $\theta$  & 0.100 & 0.084 & 15.65\% & 0.086 & 14.17\% \\
    $\sigma$  & 0.400 & 0.546 & 36.61\% & 0.714 & 78.49\% \\
    $\rho$    & $-0.600$ & $-0.550$ & 8.34\% & $-0.544$ & 9.36\% \\
    $v_0$     & 0.050 & 0.057 & 14.17\% & 0.040 & 19.93\% \\
    \midrule
    Feller satisfied?  & Yes & \multicolumn{2}{c}{Yes ($2\kappa\theta=0.644 > \sigma^2=0.299$)} & \multicolumn{2}{c}{Not enforced} \\
    \midrule
    Calibration time   & ---  & \multicolumn{2}{c}{11.76 s} & \multicolumn{2}{c}{11.30 s} \\
    Best loss          & ---  & \multicolumn{2}{c}{$3.963\times10^{-5}$ (Vega-wtd.\ IV MSE)} & \multicolumn{2}{c}{$4.045\times10^{-6}$ (normalised price MSE)} \\
    \bottomrule
  \end{tabular}
  \begin{tablenotes}
    \small
    \item \textit{Note:} Relative error $= |\hat{\theta}_i - \theta_i^*| / |\theta_i^*|$. The two methods optimise different objective functions and their best-loss values are therefore not directly comparable. The proposed method enforces the Feller condition ($2\kappa\theta > \sigma^2$) as a soft penalty during calibration; Zhang~(2025) applies no such physical constraint.
  \end{tablenotes}
\end{table}


% ============================================================
% Table 3: Re-pricing Accuracy — Proposed vs Zhang (2025)
% ============================================================
\begin{table}[htbp]
  \centering
  \caption{Re-pricing Accuracy on the Synthetic Option Chain (99 Contracts)}
  \label{tab:repricing_accuracy}
  \begin{tabular}{lcc}
    \toprule
    \textbf{Metric} & \textbf{Proposed Method} & \textbf{Zhang (2025)} \\
    \midrule
    \multicolumn{3}{l}{\textit{Implied Volatility Space}} \\
    \quad Mean IV relative error (MRE)  & 3.08\%  & 11.04\% \\
    \quad Max IV relative error         & 9.25\%  & 62.84\% \\
    \midrule
    \multicolumn{3}{l}{\textit{Price Space}} \\
    \quad Mean price relative error     & 12.78\%   & 184.80\% \\
    \quad Max price relative error      & 129.51\%  & 7873.29\% \\
    \bottomrule
  \end{tabular}
  \begin{tablenotes}
    \small
    \item \textit{Note:} The proposed method calibrates in implied volatility (IV) space with Vega weighting, and re-pricing errors are reported both in IV space and in price space (via QuantLib). Zhang~(2025) calibrates in normalised price space; its IV-space errors are computed by converting DDN-predicted prices to Black--Scholes implied volatilities. Large price errors for both methods are concentrated in deep out-of-the-money (OTM) contracts ($K/S_0 > 1.3$), where option prices are near zero and small absolute deviations produce disproportionately large relative errors. The proposed method achieves substantially lower errors in both IV and price space compared to Zhang~(2025), confirming that IV-space calibration with Feller regularisation provides more physically consistent parameter estimates.
  \end{tablenotes}
\end{table}
